\section*{Project 2: 2D FEM Analysis of a Microstrip Transmission Line}

\subsection*{Main Objective}

The aim of this project is to use the Finite Element Method (FEM) to model a two-dimensional microstrip transmission-line cross-section and calculate its characteristic impedance \(Z_0\).

The overall workflow is

\[
\text{Geometry}
\rightarrow
\text{PDE}
\rightarrow
\text{FEM}
\rightarrow
\text{Capacitance}
\rightarrow
\varepsilon_{\mathrm{eff}}
\rightarrow
Z_0.
\]

The main goal is not only to obtain a value close to \(50~\Omega\), but also to provide sufficient numerical evidence to show that the FEM result is reliable and reproducible.

\subsection*{1. Physical and Mathematical Model}

The microstrip cross-section contains:

\begin{itemize}
    \item an FR-4 substrate;
    \item an air region;
    \item a copper signal trace;
    \item a ground plane.
\end{itemize}

The nominal geometric and material parameters are

\[
H = 1.6~\mathrm{mm},
\qquad
w = 3.0~\mathrm{mm},
\qquad
t = 0.035~\mathrm{mm},
\]

with substrate relative permittivity

\[
\varepsilon_r = 4.4.
\]

The electrostatic governing equation is

\[
-\nabla \cdot \left( \varepsilon \nabla \phi \right) = 0.
\]

The main boundary conditions are:

\begin{itemize}
    \item signal trace: \(\phi = V_0\);
    \item ground plane: \(\phi = 0\);
    \item artificial outer boundaries: homogeneous Neumann boundary conditions.
\end{itemize}

The conductor interior is excluded from the field domain and treated as an internal Dirichlet boundary.

\subsection*{2. FEM Formulation}

The weak form of the problem is

\[
\int_{\Omega}
\varepsilon \nabla \phi \cdot \nabla v
\, d\Omega
=
0,
\qquad
\forall v \in V_0.
\]

The problem is discretised using conforming triangular \(P_1\) finite elements.

The FEM solution should provide:

\begin{itemize}
    \item electric potential \(\phi_h\);
    \item electric-field magnitude \(\lvert E_h \rvert\);
    \item electrostatic energy density.
\end{itemize}

The following numerical checks should also be performed:

\begin{itemize}
    \item potential bounds;
    \item symmetry about \(x=0\);
    \item matrix symmetry;
    \item positive definiteness after Dirichlet elimination;
    \item algebraic residual;
    \item global charge balance.
\end{itemize}

\subsection*{3. Solver Verification}

Before using the FEM solver for the microstrip problem, the numerical implementation should be tested using a problem with a known exact or manufactured solution.

A mesh-refinement study should be carried out and should report:

\begin{itemize}
    \item error in a specified norm;
    \item observed convergence order;
    \item algebraic residual.
\end{itemize}

This step is used to confirm that the FEM implementation is working correctly before relying on the microstrip results.

\subsection*{4. Domain-Truncation Study}

The real air region is unbounded, but the computational model must use a finite domain.

The following domain sizes should be tested:

\[
\left(
\frac{W}{H},
\frac{H_{\mathrm{air}}}{H}
\right)
=
(3,3),
(5,5),
(8,8),
(12,12).
\]

The near-trace mesh should remain fixed while the outer domain is extended.

The capacitance values should be compared using the \(12H\) case as the tested reference.

The smallest tested domain with less than \(1\%\) difference should be identified.

The largest domain should also be repeated using a finer near-field mesh to check that the domain conclusion is not caused by discretisation error.

\subsection*{5. Corner Singularity and Uniform Mesh Convergence}

The electric field becomes strongly concentrated near the conductor corners, particularly near the lower corners where the conductor meets the air--substrate interface.

A uniform mesh-refinement study should therefore be performed using the energy norm

\[
\|e\|_E^2
=
\int_{\Omega}
\varepsilon
\left|
\nabla e
\right|^2
\, d\Omega.
\]

Since no exact solution is available for the full microstrip problem, a substantially finer FEM solution can be used as the numerical reference.

When comparing solutions from different meshes, the coarse solution must be transferred correctly to the reference mesh before calculating the error.

\subsection*{6. Graded or Local Refinement Study}

A graded or locally refined mesh should also be tested, with more elements placed near the conductor corners.

The graded strategy should be compared with uniform refinement using:

\begin{itemize}
    \item error versus degrees of freedom;
    \item error versus matrix nonzeros;
    \item error versus computational time.
\end{itemize}

The aim is to determine whether local refinement provides better cost--accuracy efficiency.

The result should be reported honestly even if the tested graded strategy does not outperform uniform refinement.

\subsection*{7. Capacitance Calculation}

The electrostatic energy is

\[
W_e
=
\frac{1}{2}
\int_{\Omega}
\varepsilon
\left|
\nabla \phi
\right|^2
\, d\Omega.
\]

The energy-based capacitance is then

\[
C_E
=
\frac{2W_e}{V_0^2}.
\]

A charge-based capacitance \(C_Q\) should also be calculated from the conductor reaction.

Agreement between \(C_E\) and \(C_Q\) can be used as an implementation consistency check, but mesh convergence must still be demonstrated separately.

\subsection*{8. All-Air Solve and Effective Permittivity}

After selecting the final FEM mesh, the problem should be solved again with

\[
\varepsilon_r = 1
\]

everywhere.

The geometry, mesh, domain, element order, boundary conditions and solver settings must remain unchanged.

This gives the all-air capacitance

\[
C_{\mathrm{air}}.
\]

The effective permittivity is then calculated using

\[
\varepsilon_{\mathrm{eff}}
=
\frac{C}{C_{\mathrm{air}}}.
\]

A basic physical check is

\[
1
<
\varepsilon_{\mathrm{eff}}
<
4.4.
\]

\subsection*{9. Characteristic Impedance}

The characteristic impedance is calculated from the matched capacitance pair:

\[
Z_0
=
\frac{1}
{c_0 \sqrt{C C_{\mathrm{air}}}},
\]

where \(c_0\) is the speed of light in vacuum.

This is the final engineering result obtained from the FEM model.

\subsection*{10. Engineering Validation}

The FEM result should be compared with common microstrip engineering approximations, including:

\begin{itemize}
    \item Hammerstad--Jensen;
    \item IPC-2141.
\end{itemize}

These formulas are used as independent engineering benchmarks.

Agreement with these formulas can support the physical reasonableness of the FEM result, but it does not replace the numerical verification and mesh-convergence studies.

\subsection*{Final Deliverables}

The final project should include:

\begin{itemize}
    \item a LaTeX report of at least 8 pages;
    \item reproducible FEM source code;
    \item a README explaining how to run the project;
    \item numerical results and figures;
    \item presentation slides in PDF format;
    \item AI transparency information;
    \item individual contribution statements.
\end{itemize}

The presentation should summarise the complete numerical story:

\[
\text{Background}
\rightarrow
\text{Model}
\rightarrow
\text{Discretisation}
\rightarrow
\text{Verification}
\rightarrow
\text{Numerical Experiments}
\rightarrow
\text{Conclusions}.
\]

The presentation duration is approximately 10 minutes, followed by 3 minutes of questions.

\subsection*{Summary}

In summary, the task is not simply to calculate the microstrip characteristic impedance.

The project requires us to:

\begin{enumerate}
    \item build the physical and mathematical model;
    \item implement and verify the FEM solver;
    \item study domain truncation;
    \item investigate mesh convergence and corner effects;
    \item compare uniform and graded refinement;
    \item calculate \(C\), \(C_{\mathrm{air}}\), \(\varepsilon_{\mathrm{eff}}\) and \(Z_0\);
    \item compare the FEM result with engineering formulas;
    \item provide reproducible numerical evidence showing why the final result can be trusted.
\end{enumerate}
